% Options for packages loaded elsewhere
\PassOptionsToPackage{unicode}{hyperref}
\PassOptionsToPackage{hyphens}{url}
\documentclass[
]{article}
\usepackage{xcolor}
\usepackage[margin=1in]{geometry}
\usepackage{amsmath,amssymb}
\setcounter{secnumdepth}{-\maxdimen} % remove section numbering
\usepackage{iftex}
\ifPDFTeX
  \usepackage[T1]{fontenc}
  \usepackage[utf8]{inputenc}
  \usepackage{textcomp} % provide euro and other symbols
\else % if luatex or xetex
  \usepackage{unicode-math} % this also loads fontspec
  \defaultfontfeatures{Scale=MatchLowercase}
  \defaultfontfeatures[\rmfamily]{Ligatures=TeX,Scale=1}
\fi
\usepackage{lmodern}
\ifPDFTeX\else
  % xetex/luatex font selection
\fi
% Use upquote if available, for straight quotes in verbatim environments
\IfFileExists{upquote.sty}{\usepackage{upquote}}{}
\IfFileExists{microtype.sty}{% use microtype if available
  \usepackage[]{microtype}
  \UseMicrotypeSet[protrusion]{basicmath} % disable protrusion for tt fonts
}{}
\makeatletter
\@ifundefined{KOMAClassName}{% if non-KOMA class
  \IfFileExists{parskip.sty}{%
    \usepackage{parskip}
  }{% else
    \setlength{\parindent}{0pt}
    \setlength{\parskip}{6pt plus 2pt minus 1pt}}
}{% if KOMA class
  \KOMAoptions{parskip=half}}
\makeatother
\usepackage{color}
\usepackage{fancyvrb}
\newcommand{\VerbBar}{|}
\newcommand{\VERB}{\Verb[commandchars=\\\{\}]}
\DefineVerbatimEnvironment{Highlighting}{Verbatim}{commandchars=\\\{\}}
% Add ',fontsize=\small' for more characters per line
\usepackage{framed}
\definecolor{shadecolor}{RGB}{248,248,248}
\newenvironment{Shaded}{\begin{snugshade}}{\end{snugshade}}
\newcommand{\AlertTok}[1]{\textcolor[rgb]{0.94,0.16,0.16}{#1}}
\newcommand{\AnnotationTok}[1]{\textcolor[rgb]{0.56,0.35,0.01}{\textbf{\textit{#1}}}}
\newcommand{\AttributeTok}[1]{\textcolor[rgb]{0.13,0.29,0.53}{#1}}
\newcommand{\BaseNTok}[1]{\textcolor[rgb]{0.00,0.00,0.81}{#1}}
\newcommand{\BuiltInTok}[1]{#1}
\newcommand{\CharTok}[1]{\textcolor[rgb]{0.31,0.60,0.02}{#1}}
\newcommand{\CommentTok}[1]{\textcolor[rgb]{0.56,0.35,0.01}{\textit{#1}}}
\newcommand{\CommentVarTok}[1]{\textcolor[rgb]{0.56,0.35,0.01}{\textbf{\textit{#1}}}}
\newcommand{\ConstantTok}[1]{\textcolor[rgb]{0.56,0.35,0.01}{#1}}
\newcommand{\ControlFlowTok}[1]{\textcolor[rgb]{0.13,0.29,0.53}{\textbf{#1}}}
\newcommand{\DataTypeTok}[1]{\textcolor[rgb]{0.13,0.29,0.53}{#1}}
\newcommand{\DecValTok}[1]{\textcolor[rgb]{0.00,0.00,0.81}{#1}}
\newcommand{\DocumentationTok}[1]{\textcolor[rgb]{0.56,0.35,0.01}{\textbf{\textit{#1}}}}
\newcommand{\ErrorTok}[1]{\textcolor[rgb]{0.64,0.00,0.00}{\textbf{#1}}}
\newcommand{\ExtensionTok}[1]{#1}
\newcommand{\FloatTok}[1]{\textcolor[rgb]{0.00,0.00,0.81}{#1}}
\newcommand{\FunctionTok}[1]{\textcolor[rgb]{0.13,0.29,0.53}{\textbf{#1}}}
\newcommand{\ImportTok}[1]{#1}
\newcommand{\InformationTok}[1]{\textcolor[rgb]{0.56,0.35,0.01}{\textbf{\textit{#1}}}}
\newcommand{\KeywordTok}[1]{\textcolor[rgb]{0.13,0.29,0.53}{\textbf{#1}}}
\newcommand{\NormalTok}[1]{#1}
\newcommand{\OperatorTok}[1]{\textcolor[rgb]{0.81,0.36,0.00}{\textbf{#1}}}
\newcommand{\OtherTok}[1]{\textcolor[rgb]{0.56,0.35,0.01}{#1}}
\newcommand{\PreprocessorTok}[1]{\textcolor[rgb]{0.56,0.35,0.01}{\textit{#1}}}
\newcommand{\RegionMarkerTok}[1]{#1}
\newcommand{\SpecialCharTok}[1]{\textcolor[rgb]{0.81,0.36,0.00}{\textbf{#1}}}
\newcommand{\SpecialStringTok}[1]{\textcolor[rgb]{0.31,0.60,0.02}{#1}}
\newcommand{\StringTok}[1]{\textcolor[rgb]{0.31,0.60,0.02}{#1}}
\newcommand{\VariableTok}[1]{\textcolor[rgb]{0.00,0.00,0.00}{#1}}
\newcommand{\VerbatimStringTok}[1]{\textcolor[rgb]{0.31,0.60,0.02}{#1}}
\newcommand{\WarningTok}[1]{\textcolor[rgb]{0.56,0.35,0.01}{\textbf{\textit{#1}}}}
\usepackage{graphicx}
\makeatletter
\newsavebox\pandoc@box
\newcommand*\pandocbounded[1]{% scales image to fit in text height/width
  \sbox\pandoc@box{#1}%
  \Gscale@div\@tempa{\textheight}{\dimexpr\ht\pandoc@box+\dp\pandoc@box\relax}%
  \Gscale@div\@tempb{\linewidth}{\wd\pandoc@box}%
  \ifdim\@tempb\p@<\@tempa\p@\let\@tempa\@tempb\fi% select the smaller of both
  \ifdim\@tempa\p@<\p@\scalebox{\@tempa}{\usebox\pandoc@box}%
  \else\usebox{\pandoc@box}%
  \fi%
}
% Set default figure placement to htbp
\def\fps@figure{htbp}
\makeatother
\setlength{\emergencystretch}{3em} % prevent overfull lines
\providecommand{\tightlist}{%
  \setlength{\itemsep}{0pt}\setlength{\parskip}{0pt}}
\usepackage{bookmark}
\IfFileExists{xurl.sty}{\usepackage{xurl}}{} % add URL line breaks if available
\urlstyle{same}
\hypersetup{
  pdftitle={Correlation Heatmaps},
  hidelinks,
  pdfcreator={LaTeX via pandoc}}

\title{Correlation Heatmaps}
\author{}
\date{\vspace{-2.5em}2026-04-17}

\begin{document}
\maketitle

\subsection{Introduction}\label{introduction}

A correlation heatmap renders a correlation matrix as a grid of
colour-encoded cells, with the cell colour mapping the correlation
coefficient on a diverging palette centred at zero. It is the natural
overview figure for a multivariate dataset: dozens of pairwise
correlations that would be impossible to read as a numeric table become
a single image where blocks of similarly correlated variables stand out
by colour. Hierarchical reordering of rows and columns brings related
variables together and reveals cluster structure, often suggesting
candidate factors for downstream analysis.

\subsection{Prerequisites}\label{prerequisites}

A working understanding of correlation coefficients (Pearson, Spearman,
Kendall), basic familiarity with hierarchical clustering, and
\texttt{ggplot2} colour scales.

\subsection{Theory}\label{theory}

A correlation heatmap displays \(r_{ij}\) in cell \((i, j)\). Standard
variants:

\begin{itemize}
\tightlist
\item
  Symmetric with numbers shown --- useful when the matrix is small (≤ 8
  variables).
\item
  Upper or lower triangle only --- avoids redundancy and saves visual
  space.
\item
  Significance overlay --- marks non-significant correlations as blank
  or with a cross, so the eye is not drawn to noise.
\item
  Hierarchical reordering --- uses correlation distance (\(1 - |r|\)) to
  cluster related variables and reveal block structure.
\end{itemize}

The colour scale must be diverging (typically red--white--blue or the
Brewer \texttt{RdBu}) and centred at zero so positive and negative
correlations are visually symmetric. Sequential or qualitative palettes
destroy the sign information and are wrong for correlations.

\subsection{Assumptions}\label{assumptions}

A valid correlation matrix (any kind: Pearson for linear, Spearman for
monotone, polychoric for ordinal). The matrix should be
positive-semi-definite if it will feed downstream methods; this is
automatic when computed from raw data with \texttt{cor()}.

\subsection{R Implementation}\label{r-implementation}

\begin{Shaded}
\begin{Highlighting}[]
\FunctionTok{library}\NormalTok{(corrplot); }\FunctionTok{library}\NormalTok{(ggcorrplot)}

\NormalTok{corr\_mat }\OtherTok{\textless{}{-}} \FunctionTok{cor}\NormalTok{(mtcars)}

\CommentTok{\# corrplot}
\FunctionTok{corrplot}\NormalTok{(corr\_mat, }\AttributeTok{method =} \StringTok{"color"}\NormalTok{, }\AttributeTok{order =} \StringTok{"hclust"}\NormalTok{,}
         \AttributeTok{addCoef.col =} \StringTok{"black"}\NormalTok{, }\AttributeTok{tl.col =} \StringTok{"black"}\NormalTok{)}

\CommentTok{\# ggcorrplot with p{-}value overlay}
\NormalTok{p\_mat }\OtherTok{\textless{}{-}} \FunctionTok{cor\_pmat}\NormalTok{(mtcars)}
\FunctionTok{ggcorrplot}\NormalTok{(corr\_mat, }\AttributeTok{hc.order =} \ConstantTok{TRUE}\NormalTok{, }\AttributeTok{type =} \StringTok{"lower"}\NormalTok{,}
           \AttributeTok{p.mat =}\NormalTok{ p\_mat, }\AttributeTok{insig =} \StringTok{"blank"}\NormalTok{,}
           \AttributeTok{lab =} \ConstantTok{TRUE}\NormalTok{, }\AttributeTok{lab\_size =} \DecValTok{3}\NormalTok{)}
\end{Highlighting}
\end{Shaded}

\subsection{Output \& Results}\label{output-results}

Two correlation heatmaps for the 11-variable \texttt{mtcars} dataset: a
full symmetric heatmap with hierarchical reordering and numeric cell
labels, and a lower-triangle version with non-significant cells masked.
The first emphasises the full structure; the second is a cleaner
publication figure.

\subsection{Interpretation}\label{interpretation}

A reporting sentence (figure caption): ``Pearson correlation matrix for
11 vehicle attributes (n = 32), reordered hierarchically using
\(1 - |r|\) as the distance; cells with \(p > 0.05\) are blanked. The
two visible blocks correspond to size-related variables (weight,
displacement, horsepower) and efficiency-related variables (mpg,
quarter-mile time).'' Always state the correlation type and the
reordering method.

\subsection{Practical Tips}\label{practical-tips}

\begin{itemize}
\tightlist
\item
  Use a diverging palette centred at zero; the Brewer \texttt{RdBu} or
  \texttt{viridis::cividis} (with the centre stretched) work well in
  print.
\item
  Hierarchical reordering (\texttt{hc.order\ =\ TRUE} in
  \texttt{ggcorrplot}, \texttt{order\ =\ "hclust"} in \texttt{corrplot})
  brings related variables together and is almost always preferable to
  alphabetical order.
\item
  Mask non-significant cells (\texttt{insig\ =\ "blank"}) to direct the
  reader's attention to robust signals; alternatively, mark with crosses
  to show the test was performed.
\item
  For more than 30 variables, the figure becomes a coloured pixel grid
  without readable labels; consider summarising via principal components
  or factor analysis instead.
\item
  State the correlation type (Pearson, Spearman, Kendall) and the test
  correction (if any) in the caption; readers should not have to guess.
\item
  For partial correlations, use \texttt{ppcor::pcor()} and feed the
  resulting matrix into the same heatmap geom; partial correlations
  often reveal structure that pairwise correlations obscure.
\end{itemize}

\subsection{R Packages Used}\label{r-packages-used}

\texttt{corrplot} for the canonical correlation heatmap with reordering
and significance overlays; \texttt{ggcorrplot} for a
\texttt{ggplot2}-native interface; \texttt{Hmisc::rcorr()} for
correlation matrices with associated \(p\)-values;
\texttt{psych::corPlot()} for an alternative with rich cell formatting.

\subsection{For Reviewers}\label{for-reviewers}

\textbf{What to look for in a paper using this method.}

\begin{itemize}
\tightlist
\item
  \textbf{Common misapplications.}
\item
  \textbf{Diagnostics that should be reported but often aren't.}
\item
  \textbf{Red flags in tables and figures.}
\item
  \textbf{What to verify.}
\item
  \textbf{What an adequate Methods paragraph must contain.}
\end{itemize}

\subsection{See also --- labs in this
chapter}\label{see-also-labs-in-this-chapter}

\begin{itemize}
\tightlist
\item
  \href{labs/lab-ggplot2-grammar-and-multi-panel-layouts.Rmd}{ggplot2
  grammar and multi-panel layouts}
\end{itemize}

\end{document}
